\documentclass[11pt]{article}
\usepackage[margin=1in]{geometry}
\usepackage{amsmath,amssymb,amsthm,hyperref}
\newtheorem{theorem}{Theorem}
\newtheorem{lemma}{Lemma}
\title{Four-state thermal-overlap classification: working proof}
\author{Research working draft}
\date{11 October 2026}
\begin{document}\maketitle
\section{A finite classification for every simple four-state slow mode}\label{sec:universal}
The triplet test uses local rate information. A replacement can instead use the
full generator's slow eigenvector. The result below is necessary and sufficient
for ordinary \emph{asymptotic} cooling and heating effects throughout the working
domain: four distinct energies, a positive finite bath temperature, a reversible
irreducible generator, and a simple slow nonzero eigenvalue. It does not classify
all transient crossings or a degenerate slow eigenspace. Computing the full slow
mode is part of the test; this is not another local rate-only criterion.

\subsection{The geometry of Gibbs derivatives}
Shift the energies to $(0,a,b,c)$ with $0<a<b<c$, without changing their gaps.
Write $\pi_i(\beta)=e^{-\beta E_i}/Z(\beta)$, and let a prime always mean a
derivative in the initial inverse temperature. Define
\[
g=\pi_1'>0,\qquad u=\frac{\pi_2'}{\pi_1'},\qquad
v=\frac{\pi_3'}{\pi_1'}.
\]
For any fixed vector $w$, including the left slow mode,
\begin{equation}\label{eq:finitephi}
\frac{d}{d\beta}(w^\top\pi)=g\Phi(\beta),\qquad
\Phi=A+Bu+Cv,\quad (A,B,C)=(w_1-w_4,w_2-w_4,w_3-w_4).
\end{equation}
The identity follows from $\sum_i\pi_i'=0$.

\begin{lemma}[Strictly convex derivative curve]\label{lem:convexcurve}
For finite real $\beta$, $u'(\beta)<0$ and $v$, considered as a function
of $u$, has strictly positive second derivative. For $\beta\geq0$, put
\[
c_0=ab(b-a),\quad c_1=ac(c-a),\quad c_2=bc(c-b),\quad
\gamma=\frac{c-a}{c-b}>1,\qquad x=e^{-(c-b)\beta}.
\]
Then
\begin{equation}\label{eq:slopescalar}
\frac{dv}{du}=s(x)=\frac{c_2x^\gamma-c_0}{c_0+c_1x},\qquad
s'(x)=\frac{c_0c_2\gamma x^{\gamma-1}
+c_1c_2(\gamma-1)x^\gamma+c_0c_1}{(c_0+c_1x)^2}>0.
\end{equation}
Here $s(0)=-1$, while $s(1)>-1$ need not be positive.
\end{lemma}
\begin{proof}
For functions $h_1,\ldots,h_k$ let $W(h_1,\ldots,h_k)$ denote the
determinant with derivative orders $0,\ldots,k-1$ in its rows.
The following identity holds for any number $n$ of distinct ordered energies,
with $f_i=e^{-\beta E_i}$ and $Z=\sum_i f_i$:
\begin{equation}\label{eq:gibbswronskian}
W(\pi_1',\ldots,\pi_k')=
\frac{(-1)^{k(k-1)/2}}{Z^{k+1}}
\left(\prod_{i<j\leq k}(E_j-E_i)\right)
\left(\prod_{i=1}^k f_i\right)
\sum_{j=k+1}^n f_j\prod_{i=1}^k(E_j-E_i),\quad k<n.
\end{equation}
Indeed, adjoining a constant column gives
$W(\pi_1',\ldots,\pi_k')=W(1,\pi_1,\ldots,\pi_k)$.
Multiplying all $k+1$ functions by $Z$ multiplies their Wronskian by
$Z^{k+1}$: the product-rule transformation of the derivative rows is triangular
with diagonal $Z$. Expand the first column of $W(Z,f_1,\ldots,f_k)$.
Terms with index at most $k$ vanish by repeated columns; every remaining term
is an exponential Vandermonde determinant with the displayed sign.

For $k=2$ the identity gives $W(g,gu)=g^2u'<0$.
For $k=3$ it gives
$W(g,gu,gv)=g^3(u')^3d^2v/du^2<0$, proving strict convexity.
For completeness, a mixed two-column calculation gives
\[
W(\pi_1',\pi_3')=
\frac{f_1f_3}{Z^3}\{ab(b-a)f_2-bc(c-b)f_4\},
\]
whereas
$W(\pi_1',\pi_2')=-f_1f_2\{ab(b-a)f_3+ac(c-a)f_4\}/Z^3$.
Their ratio is $dv/du$ and simplifies to \eqref{eq:slopescalar}.
Differentiation gives the strictly positive displayed numerator.
\end{proof}

\subsection{The finite test and its equality cases}
The endpoint values needed below are explicit:
\begin{align*}
u(0)&=1-\frac{4a}{a+b+c},&v(0)&=1-\frac{4b}{a+b+c},\\
u(\infty)&=-1,&v(\infty)&=0.
\end{align*}
At the bath and any other finite temperature, stable formulas are
\[
u=\frac{-a+(b-a)e^{-b\beta}+(c-a)e^{-c\beta}}
{a+be^{-(b-a)\beta}+ce^{-(c-a)\beta}},
\]
\[
v=\frac{e^{-(b-a)\beta}\{-b+(a-b)e^{-a\beta}+(c-b)e^{-c\beta}\}}
{a+be^{-(b-a)\beta}+ce^{-(c-a)\beta}}.
\]

\begin{theorem}[Necessary and sufficient finite test]\label{thm:finiteclassification}
Let $w$ be a nonzero left eigenvector of the simple slow nonzero mode, with
$w^\top\pi(\beta_b)=0$. Form $A,B,C,\Phi$ as in \eqref{eq:finitephi}.
If $C\ne0$ and $-1<-B/C<s(1)$, define $\beta_*$ by the unique solution
\[
s(x_*)=-B/C,\quad 0<x_*<1,\qquad
\beta_*=-\frac{\log x_*}{c-b}.
\]
Otherwise omit $\beta_*$. Form the two finite sets of real numbers
\begin{align*}
\mathcal V_c&=\{\Phi(0),\Phi(\beta_b)\}
\cup\{\Phi(\beta_*):0<\beta_*<\beta_b\},\\
\mathcal V_h&=\{\Phi(\beta_b),w_1-w_2\}
\cup\{\Phi(\beta_*):\beta_b<\beta_*<\infty\}.
\end{align*}
An ordinary asymptotic cooling effect exists if and only if
$\min\mathcal V_c<0<\max\mathcal V_c$. An ordinary asymptotic heating effect
exists if and only if $\min\mathcal V_h<0<\max\mathcal V_h$.
All inequalities are strict. A zero extremum without opposite signs does not
certify an effect. Endpoints $\beta=0,\infty$ are used only as limits; every
effect asserted by the theorem has witnesses at finite positive temperatures.
\end{theorem}
\begin{proof}
If $C=0$, then $\Phi=A+Bu$ is monotone (or constant). Otherwise
$\Phi'=u'(B+C\,s(x))$. Lemma~\ref{lem:convexcurve} shows that this has at
most one zero for $\beta>0$, precisely the $\beta_*$ just specified. If
$-B/C$ equals either limiting slope, there is no interior zero. Thus the
listed values give the infimum and supremum of $\Phi$ on each side of the
bath, including its limiting endpoint. Since $g>0$, they decide whether
$b'(\beta)$ takes both positive and negative values on that open side.

It remains to justify the equivalence with a Mpemba ordering; a stationary
point alone would not suffice. Write $b=w^\top\pi$, so $b(\beta_b)=0$.
If $b'$ has only one sign on a side, $b$ is monotone there and $|b|$ increases
with distance from the bath. There is no anomalous pair. Conversely, suppose
$b'$ takes both signs. The strictly convex representation implies at most
two zeros of $b'$ in total, counting multiplicities, unless it vanishes
identically; the latter would require $w$ constant and is excluded. On the
cooling side consider the last sign-changing zero before the bath. On the
interval immediately to its right, $b$ has sign opposite to $b'$, because
$b(\beta_b)=0$ and there are no further sign changes. Immediately to its
left $b$ keeps that nonzero sign but $b'$ changes sign. Consequently
$\partial_\beta|b|>0$ there, yielding a finite cooling witness. For heating,
use the first sign-changing zero after the bath. Immediately to its left
$b$ has the same sign as $b'$; immediately to its right the derivative changes
sign while $b$ does not, so $\partial_\beta|b|<0$. This yields a finite
heating witness. Isolated tangencies do not affect this reasoning. If the
bath itself is stationary, apply the same argument on the open side.
Strict opposite signs at a limiting endpoint persist at nearby finite
temperatures by continuity.
\end{proof}

The theorem replaces an unbounded scan of initial temperatures by endpoint
evaluations and at most one scalar equation with a unique solution. Its scope
is every generator in the stated simple-mode working domain, with no restriction
on energy spacing, rate scales, or triplet branches. It is an exact mathematical
criterion, not a claim that floating-point arithmetic decides every equality
case. With rational commensurate energies in an exponential coordinate and
algebraic mode data, polynomial root isolation can certify the test. General
transcendental boundary cases require exact arguments or suitable enclosures.

\paragraph{Strong effects can also be classified with finitely many tests.}
The same construction isolates every zero of $b'$: on either side of the
optional $\beta_*$, $\Phi$ is strictly monotone unless $B=C=0$. Opposite
endpoint signs give one unique root; an exact zero at $\beta_*$ is a tangency.
There are at most two finite positive roots in total. On each side of the
bath, insert these stationary temperatures between the bath and the limiting
endpoint ($0$ or $\infty$). Evaluate $b$ at the resulting ordered knots.
Every pair of consecutive nonzero values of opposite sign gives exactly one
strong preparation in that open interval, since $b$ is strictly monotone
there. Every stationary knot with $b=0$, other than the bath, is itself a
strong preparation, including a tangential zero. A zero at $0$ or $\infty$
is excluded. Intervals adjacent to a zero knot contain no additional zero.
This is a complete finite classification of nonbath strong preparations;
their actual temperatures, if desired, are the unique roots on those monotone
intervals. It preserves the distinction between an ordinary turning point and
a strong zero. The numerical implementation supplied here currently implements
the ordinary-effect test; this strong-effect extension is a mathematical
root-isolation prescription.

\subsection{A sharp limit on thermal complexity}
For four states, Lemma~\ref{lem:convexcurve} shows that $b'$ has at most two
zeros counting multiplicity: a line intersects a strictly convex graph at
most twice, with tangency counting twice. More generally,
\eqref{eq:gibbswronskian} shows that $(\pi_1',\ldots,\pi_{n-1}')$ is an
extended complete Chebyshev system. This means every nonzero linear combination
of its first $k$ functions has at most $k-1$ zeros counting multiplicity.
The implication from nonvanishing prefix Wronskians is the standard repeated
weighted Rolle argument (see, for example,
\href{https://repository.tudelft.nl/file/File_ce2c559d-93ee-46e3-8128-227b977215ac}{Gasull, Geyer and Ma\~nosas,
\emph{A Chebyshev criterion with applications}, J. Differential Equations 269
(2020), Lemma 2.3}): divide by the first nonvanishing function,
differentiate, and repeat; the transformed prefix Wronskians remain nonzero.
At each differentiation the number of zeros can decrease by at most one.
Thus any nonconstant Gibbs expectation has at most $n-2$ stationary points.
This is a consequence of established Chebyshev-system theory; no new general
Wronskian or zero-counting principle is claimed.

Likewise, $Zb=\sum_iw_i e^{-\beta E_i}$ has at most $n-1$ zeros counting
multiplicity by the exponential Vandermonde argument. One is the bath, so
there are at most $n-2$ nonbath strong preparations. These two bounds are
sharp even for positive reversible generators with a simple slow mode.
Indeed, prescribe $n-2$ distinct positive nonbath inverse temperatures.
Together with the bath they impose $n-1$ homogeneous interpolation conditions
on $n$ coefficients $w_i$. The exponential Chebyshev property makes the
nullspace one-dimensional and its nonzero solution has exactly those simple
zeros and no others. Rolle's theorem forces $n-2$ stationary points between
them, attaining the derivative bound as well.

To realize this $w$, let $D=\operatorname{diag}(\pi_b)$,
$N=w^\top Dw>0$, and choose
\begin{equation}\label{eq:rankonerealization}
M=-I+\pi_b\mathbf1^\top+\epsilon\frac{Dw w^\top}{N},\qquad
0<\epsilon<\min\left\{1,\frac{N}{\max_i|w_i|^2}\right\}.
\end{equation}
Its off-diagonal entries are $\pi_{bi}(1+\epsilon w_iw_j/N)>0$;
its column sums vanish and detailed balance holds. The eigenvalues are
$0,-1+\epsilon,-1,\ldots,-1$, with $w$ the left slow mode. This elementary
spectral realization is used to show attainability, without claiming a new
general generator-construction method. It need not obey the triplet premise.

In particular, four is the minimum number of distinct states allowing both
ordinary cooling and ordinary heating effects at one fixed bath with a simple
slow mode: each side needs a sign-changing stationary point, while three states
allow at most one globally. Four states attain the bound. Take
$E_i=(i-1)\log2$, $\beta_b=1$, and
\[
w=(-3/32,\ 11/16,\ -3/2,\ 1),\qquad
\sum_{i=1}^4w_i x^{i-1}=(x-1/4)(x-1/2)(x-3/4).
\]
The two nonbath strong preparations are
$\beta=\log(4/3)/\log2<1$ and $\beta=2>1$. Choosing $\epsilon=1/100$
in \eqref{eq:rankonerealization} gives an entirely rational generator, with
eigenvalues $0,-99/100,-1,-1$. Both effects follow from the two exact nonbath
zeros. This example illustrates attainability, not a second triplet
counterexample. Publication priority for these Gibbs-specific consequences
requires a separate literature assessment.

\end{document}
